% !PRINTED 379
\weberpage{379}
%<<CONT
\weq{14}{\begin{aligned}
a &= a_1\delta^2 + b_1\beta\delta - c_1\beta^2\\
b &= 2a_1\gamma\delta + b_1(\alpha\delta + \beta\gamma) - 2c_1\alpha\beta,\\
-c &= a_1\gamma^2 + b_1\alpha\gamma - c_1\alpha^2,
\end{aligned}}
and further
\weq{15}{b^2 + 4ac = b_1^2 + 4a_1c_1 = D.}
Since $a$, $b$, $c$ are without common divisor, $a_1$, $b_1$, $c_1$ also have no common divisor, as one sees from (14), and $a_1$, $b_1$, $c_1$ have the same meaning for $\omega_1$ as $a$, $b$, $c$ have for $\omega$.

\gesperrt{VII. One sees further from (15) that for equivalent numbers not only the irrationality $\sqrt{d}$, but also the discriminant is the same.}

If one sets
\weq{16}{\omega_1 = x_1 + y_1\sqrt{d} = \frac{\alpha(x+y\sqrt{d})+\beta}{\gamma(x+y\sqrt{d})+\delta},}
and, in order to determine $x_1$, $y_1$, multiplies the numerator and denominator of the last fraction by $\gamma(x-y\sqrt{d})+\delta$, then by (4) and (13) there results
\weq{17}{2x_1 = \frac{b_1}{c_1}, \qquad y_1 = \frac{\varepsilon y c}{c_1},}
hence, if the $\sqrt{D}$ is understood in the same way as in (7) and (8), by (9):
\weq{18}{\omega_1 = \frac{b_1 + \varepsilon\sqrt{D}}{2c_1}.}

\begin{center}
\S.\ 123.

\gesperrt{Reduced Numbers with Negative Discriminant.} % [sic] section number \S.\,123 repeats the number already used earlier on this same page for ``Reducirte Zahlen'' as printed
\end{center}

The considerations up to this point are applicable uniformly to the two cases of real and imaginary quadratic irrational numbers, that is, to positive and negative discriminants. Now, however, the two cases must be separated from one another, and we begin with the negative discriminants. We are concerned, then, with imaginary numbers:
\weq{1}{\omega = x + y\sqrt{d} = \xi + i\eta,}
in which $\xi$ and $\eta$ are real; it suffices to consider one of the two conjugate numbers $\xi+i\eta$, $\xi-i\eta$. We therefore assume $\eta$ to be positive and lay down the following definition:
%CONT>>
